% Generated by book/tools/notebook_to_tex.py. Do not edit by hand.

\chapter[BERT 첫 만남 (Pipeline)]{BERT 첫 만남 (Pipeline)}

\markboth{BERT 첫 만남 (Pipeline)}{BERT 첫 만남 (Pipeline)}

\chaptermeta{07_bert_pipeline/07_bert_pipeline.ipynb}{https://colab.research.google.com/github/yoon-gu/neuqes-101/blob/master/07_bert_pipeline/07_bert_pipeline.ipynb}{pipeline 한 줄 뒤의 tokenizer, model, post-processing 분해}



\index{1BERT@BERT}

\index{1DistilBERT@DistilBERT}

\index{1pipeline@pipeline}

\index{1AutoTokenizer@AutoTokenizer}

\index{1AutoModelForSequenceClassification@AutoModelForSequenceClassification}

\index{1WordPiece@WordPiece}

\index{1pretrained model@pretrained model}

\index{1special token@special token}

\index{1post processing@post-processing}

\index{0BERT 첫 만남@BERT 첫 만남}

\index{0사전학습 모델@사전학습 모델}

\index{0파이프라인@파이프라인}

\index{0워드피스@워드피스}

\index{0특수 토큰@특수 토큰}

\index{0토크나이저@토크나이저}

\index{0추론@추론}

\index{1transformers@transformers}

\index{1pipeline sentiment analysis@pipeline("sentiment-analysis")}

\index{1from pretrained@from\_pretrained}

\index{1model forward@model.forward}

\index{1logits@logits}

\index{1softmax@softmax}

\index{1argmax@argmax}

\index{1CLS@[CLS]}

\index{1SEP@[SEP]}

\index{1WordPiece prefix@WordPiece prefix}

\index{0WordPiece 접두사@WordPiece 접두사}

\index{1SST 2@SST-2}

\index{1sentiment analysis@sentiment analysis}

\index{0감성 분석@감성 분석}

\index{0모델 다운로드@모델 다운로드}

\index{0캐시@캐시}

\index{0토큰 ID@토큰 ID}

\index{0후처리@후처리}



\section{변화추적표}

\textbf{Phase 1 시작} --- scikit-learn 단계 끝, \inlinecode{transformers} 등장.

\begin{booktable}{7장 변화추적표}{tab:ch07-01}
\begin{adjustbox}{max width=\textwidth}
\begin{tabular}{@{}lllllll@{}}
\toprule
Ch & 모델 & 토크나이저 & 데이터 & Output Head & Activation & Loss \\
\midrule

1 & (TF-IDF) & \inlinecode{TfidfVectorizer()} & Yelp 5,000 & --- & --- & --- \\
2 & \inlinecode{LinearRegression()} & \inlinecode{TfidfVectorizer()} & Yelp (별점 1-5) & (1차원) & 없음 & \inlinecode{MSELoss} \\
3 & \inlinecode{LogisticRegression()} & \inlinecode{TfidfVectorizer()} & Yelp 이진화 & (1차원) & sigmoid & \inlinecode{BCEWithLogitsLoss} \\
4 & \inlinecode{LogisticRegression(multi\_class="multinomial")} & \inlinecode{TfidfVectorizer()} & Yelp 이진화 & (2차원) & softmax & \inlinecode{CrossEntropyLoss} \\
5 & \inlinecode{LogisticRegression(multi\_class="multinomial")} & \inlinecode{TfidfVectorizer()} & Yelp 5클래스 & (5차원) & softmax & \inlinecode{CrossEntropyLoss} \\
6 & \inlinecode{OneVsRestClassifier(LogisticRegression())} & \inlinecode{TfidfVectorizer()} & Yelp + 측면 합성 & (5차원) & per-label sigmoid & per-label \inlinecode{BCEWithLogitsLoss} \\
\textbf{7 ← 여기} & \inlinecode{pipeline("sentiment-analysis")} & \inlinecode{AutoTokenizer.from\_pretrained(...)} & 간단 영어 예시 & \textbf{사전학습 헤드} & softmax & --- (추론만) \\
\bottomrule
\end{tabular}
\end{adjustbox}
\end{booktable}

전체 20장 표는 \href{https://github.com/yoon-gu/neuqes-101\#챕터별-변화추적표}{루트 README.md}를 참고하세요.



\section{변경점: 6장 대비}

\begin{booktable}{7장 변경점 요약}{tab:ch07-02}
\begin{adjustbox}{max width=\textwidth}
\begin{tabular}{@{}lll@{}}
\toprule
축 & 6장 & 7장 \\
\midrule

라이브러리 & \inlinecode{sklearn} & \textbf{\inlinecode{transformers}} (Hugging Face) \\
모델 & \inlinecode{OneVsRestClassifier(LogisticRegression())} (학습) & \textbf{\inlinecode{pipeline("sentiment-analysis")}} (사전학습 + 추론) \\
토크나이저 & \inlinecode{TfidfVectorizer()} (단어 단위 어휘 학습) & \textbf{\inlinecode{AutoTokenizer}} (사전학습된 WordPiece) \\
학습 단계 & sklearn \inlinecode{fit()} 한 번에 학습 & \textbf{학습 없음} --- 사전학습 가중치 로드 후 추론만 \\
데이터 & Yelp 5,000건 & 간단 영어 예시 문장 (분해 시연용) \\
하드웨어 & CPU & CPU 또는 T4 GPU (이 장는 추론만이라 어느 쪽도 OK) \\
\bottomrule
\end{tabular}
\end{adjustbox}
\end{booktable}

\textbf{왜 학습 없이 시작하나?} Phase 1 첫 장은 \inlinecode{transformers} 의 \emph{추상화 계층} 을 익히는 데 집중합니다. \inlinecode{pipeline} 한 줄 뒤에 토크나이저·모델·후처리 3단계가 어떻게 굴러가는지 손에 잡히면, 8장(Tokenizer/Datasets 해부)와 9장(BERT 회귀 첫 학습)에서 \inlinecode{Trainer} 가 등장할 때 코드를 \emph{읽을} 수 있습니다.



\section{토크나이저 노트 --- 첫 WordPiece 등장}

이 장의 토크나이저는 \textbf{사전학습된 WordPiece}. Phase 0의 \inlinecode{TfidfVectorizer} 와 \emph{완전히 다른 패러다임} 입니다.

\begin{booktable}{7장 토크나이저 노트 - 첫 WordPiece 등장 표}{tab:ch07-03}
\begin{adjustbox}{max width=\textwidth}
\begin{tabular}{@{}lll@{}}
\toprule
비교 & TF-IDF (Phase 0) & WordPiece (Phase 1+) \\
\midrule

분리 단위 & 단어 (whitespace + 정규식) & \textbf{서브워드} (자주 등장하는 문자 시퀀스) \\
어휘 출처 & 학습 데이터에서 그때그때 학습 & \textbf{사전학습된 30,522개 어휘} (BERT 학습 시 정해짐) \\
OOV 처리 & 단순히 무시 & \texttt{{[}UNK{]}} 또는 더 작은 서브워드로 분해 \\
특수 토큰 & 없음 & \texttt{{[}CLS{]}}, \texttt{{[}SEP{]}}, \texttt{{[}PAD{]}}, \texttt{{[}MASK{]}} 등 \\
출력 & sparse vector (V차원, 거의 0) & 정수 ID 시퀀스 + attention mask \\
\bottomrule
\end{tabular}
\end{adjustbox}
\end{booktable}

같은 문장 \inlinecode{"I\ love\ using\ Hugging\ Face!"} 가 어떻게 토큰화되는지 곧 직접 확인합니다 (Step 2). \inlinecode{\#\#} 접두사가 보이는 단어는 어디고, 왜 그렇게 쪼개졌는지도 같이 봅니다.

\begin{previewBox}{미리보기}
\textbf{다음 장(8장)}: 같은 WordPiece 토크나이저를 \emph{깊게} --- \inlinecode{padding}, \inlinecode{truncation}, \inlinecode{max\_length} 옵션과 \inlinecode{datasets} 라이브러리 메모리 효율까지.
\end{previewBox}



\section{0. 환경 준비}

Colab에는 \inlinecode{transformers}가 보통 설치돼 있지만, 최신 버전을 보장하기 위해 한 번 설치합니다.



\Needspace{8\baselineskip}
\begin{lstlisting}
!pip install -q transformers
\end{lstlisting}



\Needspace{17\baselineskip}
\begin{lstlisting}
import torch
print(f"PyTorch:        {torch.__version__}")
print(f"CUDA available: {torch.cuda.is_available()}")
if torch.cuda.is_available():
    print(
        f"GPU:            "
        f"{torch.cuda.get_device_name(0)}"
    )
else:
    print(
        "Running on CPU (inference works in this chapter;"
        "skip the nvidia-smi cells)"
    )
\end{lstlisting}



\subsection{\texorpdfstring{\inlinecode{!nvidia-smi} --- GPU 메모리(VRAM) 실시간 추적}{!nvidia-smi --- GPU 메모리(VRAM) 실시간 추적}}

이 장부터 학습·추론 코드가 GPU에 모델을 올리기 시작합니다. \textbf{\inlinecode{!nvidia-smi}} 는 NVIDIA에서 제공하는 명령행 도구로, 현재 GPU의 VRAM 사용량·온도·전력을 한 번에 보여줍니다. Colab 셀에서 \inlinecode{!} 접두사로 호출 가능.

T4의 총 VRAM은 \textbf{약 15.36 GB} (= 15,360 MiB). 모델·옵티마이저·activation을 모두 이 안에 담아야 합니다 --- 9장 이후 학습 chapter에서는 이 한도와 자주 부딪히게 되어요.

\textbf{baseline} --- 아직 아무 모델도 안 올린 상태:



\Needspace{8\baselineskip}
\begin{lstlisting}
!nvidia-smi
\end{lstlisting}



\textbf{무엇을 봐야 하나} --- 출력 가운데 줄 \inlinecode{Memory-Usage} 칸:

\Needspace{6\baselineskip}
\begin{lstlisting}
| ... |  XXX MiB / 15360MiB | ...
        └─ used    └─ total
\end{lstlisting}

\begin{itemize}
\tightlist
\item
  처음엔 \textasciitilde3-200 MiB 정도. CUDA 컨텍스트가 잡혀 있는 만큼만.
\item
  모델을 GPU에 올릴 때마다 \inlinecode{used} 가 증가합니다.
\item
  \inlinecode{Volatile\ GPU-Util} 은 \emph{현재} GPU가 일하는 비율 --- 학습 중에는 90\textasciitilde100\% 가까이.
\end{itemize}

\textbf{Python으로도 확인 가능} (셀 내부에서 변수로 받고 싶을 때):

\Needspace{8\baselineskip}
\begin{lstlisting}
if torch.cuda.is_available():
    used  = torch.cuda.memory_allocated() / 1024**2
    total = torch.cuda.get_device_properties(0).total_memory / 1024**2
    print(f"GPU memory: {used:.0f} / {total:.0f} MiB")
\end{lstlisting}

\begin{quote}
Tip: \inlinecode{!nvidia-smi} 는 \emph{시스템 전체} VRAM을 보여주고, \inlinecode{torch.cuda.memory\_allocated()} 는 \emph{현재 PyTorch 프로세스} 의 할당량만 보여줍니다 --- 후자는 캐시·예약 메모리는 빼고 실제 텐서가 점유한 양에 가깝습니다.
\end{quote}



\section{1. 실습: 일단 돌려봅시다}

Hugging Face의 \inlinecode{pipeline} 은 \textbf{“모델 다운로드 → 토큰화 → 추론 → 결과 후처리”} 를 한 줄로 묶어주는 함수입니다.

감성 분석(sentiment analysis)부터 시작합니다. \textbf{GPU가 있으면 \inlinecode{device=0} 으로 명시} --- 그래야 모델이 VRAM에 올라가서 nvidia-smi 변화가 보입니다 (기본은 CPU).



\Needspace{10\baselineskip}
\begin{lstlisting}
from transformers import pipeline

DEVICE = 0 if torch.cuda.is_available() else -1
# 0 = GPU index, -1 = CPU
classifier = pipeline("sentiment-analysis", device=DEVICE)
classifier("I love using Hugging Face! It's so simple.")
\end{lstlisting}

\begin{codeRead}
\inlinecode{transformers} 같은 주요 패키지를 불러옵니다. \textbf{3행}의 \inlinecode{DEVICE = ...}에서는 이후 분석에서 사용할 값을 준비합니다. \textbf{5행}의 \inlinecode{classifier = ...}에서는 이후 분석에서 사용할 값을 준비합니다. \textbf{6행}의 \inlinecode{classifier(...)}에서는 앞 단계에서 만든 값을 바탕으로 다음 계산을 수행합니다.
\end{codeRead}



\textbf{DistilBERT(SST-2)가 VRAM에 올라간 직후의 nvidia-smi:}



\Needspace{8\baselineskip}
\begin{lstlisting}
!nvidia-smi
\end{lstlisting}

\begin{codeRead}
\textbf{1행}의 \inlinecode{!nvidia-smi}에서는 앞 단계에서 만든 값을 바탕으로 다음 계산을 수행합니다.
\end{codeRead}



baseline과 비교하면 \inlinecode{Memory-Usage} 가 늘어났을 겁니다. \textbf{얼마나 늘어났을까?} --- 모델 파라미터 수에서 거꾸로 추정해봅니다.



\Needspace{17\baselineskip}
\begin{lstlisting}
def model_size_summary(model, dtype_bytes=4):

    n_params = sum(p.numel() for p in model.parameters())
    size_mb = n_params * dtype_bytes / 1024**2
    return n_params, size_mb

n, size_mb = model_size_summary(classifier.model)
print(f"DistilBERT (SST-2) parameters:")
print(f"  count:           {n:>13,}  ({n/1e6:.1f} M)")
print(
    f"  fp32 size:       "
    f"{size_mb:>10.1f} MB  (= params × 4 bytes)"
)
\end{lstlisting}

\begin{codeRead}
\textbf{1행}의 \inlinecode{def model\_size\_summary(mo...}에서는 이후 분석에서 사용할 값을 준비합니다. \textbf{3행}의 \inlinecode{n\_params = ...}에서는 이후 분석에서 사용할 값을 준비합니다. \textbf{4행}의 \inlinecode{size\_mb = ...}에서는 이후 분석에서 사용할 값을 준비합니다. \textbf{5행}의 \inlinecode{return n\_params, size\_mb}에서는 앞 단계에서 만든 값을 바탕으로 다음 계산을 수행합니다. \textbf{7행}의 \inlinecode{n, size\_mb = model\_size\_s...}에서는 이후 분석에서 사용할 값을 준비합니다.
\end{codeRead}



\textbf{파라미터 수와 VRAM의 관계}

가중치 한 개는 \emph{한 개의 부동소수} --- fp32면 4 bytes, fp16/bf16이면 2 bytes를 차지합니다. 그래서:

\begin{equation}
\label{eq:ch07-01}
\text{모델 가중치 크기} \approx \text{파라미터 수} \times \text{dtype bytes}
\end{equation}

식~\eqref{eq:ch07-01}은 이 절에서 사용하는 핵심 관계를 정리한 것입니다.

\begin{booktable}{7장 1. 실습: 일단 돌려봅시다 표}{tab:ch07-04}
\begin{adjustbox}{max width=\textwidth}
\begin{tabular}{@{}lll@{}}
\toprule
dtype & bytes/param & DistilBERT(67M) 예상 크기 \\
\midrule

\textbf{fp32} (기본 학습) & 4 & \textasciitilde255 MB \\
\textbf{fp16 / bf16} (mixed precision) & 2 & \textasciitilde128 MB \\
\textbf{int8} (양자화 추론) & 1 & \textasciitilde64 MB \\
\bottomrule
\end{tabular}
\end{adjustbox}
\end{booktable}

\textbf{그런데 nvidia-smi는 파라미터 크기보다 \emph{더 많이} 늘어나는데요?} 차이의 정체:

\begin{itemize}
\tightlist
\item
  \textbf{PyTorch CUDA 컨텍스트} (\textasciitilde150-300 MB): 라이브러리·드라이버가 GPU 점유 시 한 번 잡는 고정 비용. 첫 모델 로드에서만 보이고 이후엔 누적되지 않음.
\item
  \textbf{CUDA 캐시 할당자} (수십 MB): PyTorch가 자주 쓰일 텐서를 캐싱.
\item
  \textbf{추론 중 일시 activation}: forward pass에서 layer 사이 중간 결과. 추론은 작지만 학습은 큼.
\end{itemize}

\textbf{학습이 되면 메모리는 더 커집니다} --- Adam 옵티마이저는 모델당 \emph{추가로 2배} (1차·2차 모멘텀)를 더 들고, gradient도 \emph{모델 크기만큼} 한 벌 --- 즉 학습 중엔 \textbf{fp32 기준 파라미터 × 4배 정도} 의 VRAM이 필요합니다. 9장에서 다시 다룹니다.



\textbf{참고}: 처음 실행 시 모델 다운로드(약 250MB)에 30초\textasciitilde1분 정도 걸립니다. 두 번째부터는 캐시되어 즉시 실행.

여러 문장도 한 번에:



\Needspace{11\baselineskip}
\begin{lstlisting}
results = classifier([
    "This movie was fantastic.",
    "Worst experience ever.",
    "It was okay, nothing special.",
])
for r in results:
    print(r)
\end{lstlisting}

\begin{codeRead}
\textbf{1--5행}의 \inlinecode{results = ...}에서는 이후 분석에서 사용할 값을 준비합니다. \textbf{6행}의 \inlinecode{for r in results:}에서는 앞 단계에서 만든 값을 바탕으로 다음 계산을 수행합니다.
\end{codeRead}



\subsection{다른 task도 같은 패턴}

\inlinecode{pipeline} 의 첫 인자만 바꾸면 다른 NLP 작업을 즉시 할 수 있습니다.



\Needspace{14\baselineskip}
\begin{lstlisting}
generator = pipeline(
    "text-generation",
    model="gpt2",
    device=DEVICE,
)
generator(
    "Hugging Face is",
    max_length=30,
    num_return_sequences=1,
)
\end{lstlisting}



\Needspace{10\baselineskip}
\begin{lstlisting}
unmasker = pipeline(
    "fill-mask",
    model="bert-base-uncased",
    device=DEVICE,
)
unmasker("Hugging Face is a [MASK] for NLP.")
\end{lstlisting}



\textbf{3개 pipeline(DistilBERT + GPT-2 + BERT-base)이 모두 VRAM에 쌓인 상태:}



\Needspace{8\baselineskip}
\begin{lstlisting}
!nvidia-smi
\end{lstlisting}



\textbf{파라미터 수 합계로 예측한 모델 가중치 vs 실제 VRAM 증가} --- 차이가 PyTorch 오버헤드입니다.



\Needspace{26\baselineskip}
\begin{lstlisting}
models_info = {
    "DistilBERT (SST-2)":       classifier.model,
    "GPT-2":                    generator.model,
    "BERT-base-uncased":        unmasker.model,
}

print(f"{'model':>30}  {'params':>12}  {'fp32 size':>14}")
print("-" * 65)
total_params, total_size = 0, 0.0
for name, m in models_info.items():
    n, sz = model_size_summary(m)
    total_params += n
    total_size += sz
    print(f"{name:>30}  {n/1e6:>9.1f} M  {sz:>11.1f} MB")
print("-" * 65)
print(
    f"{'total':>30}  {total_params/1e6:>9.1f} M  "
    f"{total_size:>11.1f} MB"
)
print(
    f"{'in GB':>30}  {' '*12}  "
    f"{total_size/1024:>11.2f} GB"
)
\end{lstlisting}



\textbf{실제 nvidia-smi VRAM 사용량과 비교} --- 위 합계(\textasciitilde 수백 MB)에 \textbf{PyTorch CUDA 컨텍스트 + 캐시 할당자 \textasciitilde{} 200-400 MB} 를 더한 값이 nvidia-smi 의 \inlinecode{Memory-Usage} 와 비슷할 겁니다.

\begin{quote}
모델별 파라미터 차이가 흥미롭습니다.

\begin{itemize}
\tightlist
\item
  \textbf{DistilBERT (\textasciitilde67M)}: BERT-base에서 layer를 절반으로 줄여 학습한 경량화 모델. 추론 속도 \textasciitilde2배.
\item
  \textbf{GPT-2 small (\textasciitilde124M)}: 파라미터는 BERT-base보다 약간 많고, 어휘(50K)도 더 큼.
\item
  \textbf{BERT-base (\textasciitilde110M)}: BERT 표준 사이즈.
\end{itemize}
\end{quote}

\textbf{메모리를 비우는 표준 패턴} --- 더 이상 안 쓰는 모델은:

\Needspace{9\baselineskip}
\begin{lstlisting}
import gc, torch
del generator, unmasker
gc.collect()
if torch.cuda.is_available():
    torch.cuda.empty_cache()
\end{lstlisting}

T4 메모리(15.36 GB)는 작은 추론 모델 여러 개를 무리 없이 담지만, BERT-large(\textasciitilde340M, fp32 \textasciitilde1.4 GB)나 학습 시 \emph{옵티마이저 + gradient} 까지 올리면 빠르게 한도에 도달하니 항상 nvidia-smi로 잔여 VRAM 확인하는 습관이 좋습니다.



\section{2. 해부: pipeline 안에서는 뭐가 일어났을까?}

\inlinecode{pipeline("sentiment-analysis")} 한 줄이 사실은 \textbf{3단계} 로 구성됩니다.

\Needspace{11\baselineskip}
\begin{lstlisting}
입력 텍스트
   ↓ [1] Tokenizer  (텍스트 → 숫자 ID)
input_ids, attention_mask
   ↓ [2] Model      (숫자 → 로짓)
logits
   ↓ [3] Post-processing (로짓 → 라벨)
{'label': 'POSITIVE', 'score': 0.9998}
\end{lstlisting}

\subsection{등장 인물 정리}

\begin{booktable}{7장 등장 인물 정리 표}{tab:ch07-05}
\begin{adjustbox}{max width=\textwidth}
\begin{tabular}{@{}lll@{}}
\toprule
컴포넌트 & 라이브러리 & 역할 \\
\midrule

\inlinecode{pipeline} & \inlinecode{transformers} & 위 3단계를 묶은 wrapper \\
Tokenizer & \inlinecode{transformers} (내부적으로 \inlinecode{tokenizers}) & 텍스트를 모델이 먹을 수 있는 숫자로 변환 \\
Model & \inlinecode{transformers} + \inlinecode{torch} & 실제 신경망 forward 연산 \\
\bottomrule
\end{tabular}
\end{adjustbox}
\end{booktable}

현재 \inlinecode{classifier} 객체가 어떤 모델/토크나이저를 사용하는지 확인합니다.



\Needspace{20\baselineskip}
\begin{lstlisting}
print(
    f"Model:               "
    f"{classifier.model.config._name_or_path}"
)
print(
    f"Model class:         "
    f"{type(classifier.model).__name__}"
)
print(
    f"Tokenizer class:     "
    f"{type(classifier.tokenizer).__name__}"
)
print(
    f"Label mapping:       "
    f"{classifier.model.config.id2label}"
)
\end{lstlisting}



기본 모델은 \textbf{\inlinecode{distilbert-base-uncased-finetuned-sst-2-english}} 입니다.

\begin{itemize}
\tightlist
\item
  \textbf{DistilBERT}: BERT를 40\% 작게 만든 경량화 모델 (학생 모델, 지식 증류로 만듦).
\item
  \textbf{SST-2}: 영화 리뷰 감성 분류 데이터셋 (Stanford Sentiment Treebank).
\item
  즉, “BERT를 영화 리뷰 데이터로 파인튜닝한 모델”입니다 --- 우리가 9장에서 직접 할 작업의 \emph{완성된} 버전.
\end{itemize}



\section{3. 변형: pipeline 없이 직접 해보기}

이제 \inlinecode{pipeline} 이 감춰뒀던 단계를 \textbf{직접 한 줄씩 실행} 합니다. 이 부분을 이해하면 앞으로 모든 모델을 자유롭게 다룰 수 있습니다.

\subsection{Step 1: Tokenizer와 Model 직접 로드}



\Needspace{21\baselineskip}
\begin{lstlisting}
from transformers import AutoTokenizer, AutoModelForSequenceClassification

model_name = "distilbert-base-uncased-finetuned-sst-2-english"

tokenizer = AutoTokenizer.from_pretrained(model_name)
model = AutoModelForSequenceClassification.from_pretrained(model_name)

if torch.cuda.is_available():
    model = model.to("cuda")

print("Loaded")
print(f"  tokenizer class: {type(tokenizer).__name__}")
print(f"  model class:     {type(model).__name__}")
print(
    f"  model device:    "
    f"{next(model.parameters()).device}"
)
\end{lstlisting}



\textbf{\inlinecode{pipeline} 위에 추가로 \emph{같은 DistilBERT} 가 올라간 상태} --- VRAM이 또 한 번 늘어났습니다. 같은 가중치라도 별도 객체이면 별도 메모리.



\Needspace{8\baselineskip}
\begin{lstlisting}
!nvidia-smi
\end{lstlisting}



\begin{quote}
\textbf{잠깐, \inlinecode{Auto}가 뭔가요?}

\inlinecode{AutoTokenizer}, \inlinecode{AutoModel...} 같은 \inlinecode{Auto} 계열 클래스는 모델 이름만 주면 \textbf{알아서 적합한 클래스를 골라주는 팩토리} 입니다.

\begin{itemize}
\tightlist
\item
  DistilBERT 모델 → \inlinecode{DistilBertTokenizer}, \inlinecode{DistilBertForSequenceClassification}
\item
  BERT 모델 → \inlinecode{BertTokenizer}, \inlinecode{BertForSequenceClassification}
\item
  GPT-2 모델 → \inlinecode{GPT2Tokenizer}, \inlinecode{GPT2LMHeadModel}
\end{itemize}

직접 \inlinecode{BertTokenizer.from\_pretrained(...)}라고 써도 되지만, \inlinecode{AutoTokenizer} 를 쓰면 모델만 바꿔도 코드가 그대로 동작합니다.
\end{quote}



\subsection{Step 2: 텍스트 → 숫자 (Tokenization)}



\Needspace{13\baselineskip}
\begin{lstlisting}
text = "I love using Hugging Face!"

tokens = tokenizer.tokenize(text)
print(f"Tokens: {tokens}")

inputs = tokenizer(text, return_tensors="pt")
print("\nModel inputs:")
for key, value in inputs.items():
    print(f"  {key}: {value}")
\end{lstlisting}



\textbf{관찰 포인트}

\begin{itemize}
\tightlist
\item
  \inlinecode{input\_ids}: 각 토큰의 정수 ID. 맨 앞 \inlinecode{101}은 \texttt{{[}CLS{]}}, 맨 뒤 \inlinecode{102}는 \texttt{{[}SEP{]}} 특수 토큰.
\item
  \inlinecode{attention\_mask}: \inlinecode{1}이면 “이 위치는 진짜 토큰”, \inlinecode{0}이면 “패딩이니 무시하라”는 신호 (이번엔 패딩 없음 → 모두 1).
\item
  \inlinecode{tokenize()} 와 \inlinecode{tokenizer()} 의 차이 --- 전자는 토큰 문자열만 반환, 후자는 모델 입력용 텐서까지 다 만들어줌.
\end{itemize}



\Needspace{9\baselineskip}
\begin{lstlisting}
print(f"{'ID':>5}    token")
print("-" * 30)
for token_id in inputs["input_ids"][0]:
    token = tokenizer.decode([token_id])
    print(f"{token_id.item():>5}    {token!r}")
\end{lstlisting}



\subsection{Step 3: 숫자 → 로짓 (Model forward)}



\Needspace{16\baselineskip}
\begin{lstlisting}
inputs_on_device = {k: v.to(model.device) for k, v in inputs.items()}

with torch.no_grad():
    outputs = model(**inputs_on_device)

print(f"Output object: {type(outputs).__name__}")
print(
    f"Logits shape:  "
    f"{outputs.logits.shape}  (batch=1, classes=2: NEGATIVE, POSITIVE)"
)
print(f"Logits:        {outputs.logits}")
print(f"Logits device: {outputs.logits.device}")
\end{lstlisting}



로짓(logits)은 모델이 출력한 \textbf{정규화되지 않은 점수}. shape \texttt{{[}1,\ 2{]}}는 “배치 1개, 클래스 2개”를 의미합니다.

여기서 잠깐 --- 익숙하지 않나요? 4장의 \emph{softmax + 2차원 head} 구조와 정확히 같습니다. BERT는 사전학습된 \emph{심층} 모델일 뿐, 마지막 분류 헤드는 sklearn에서 본 것과 본질이 같습니다.



\subsection{Step 4: 로짓 → 확률/라벨 (Post-processing)}



\Needspace{15\baselineskip}
\begin{lstlisting}
probs = torch.softmax(outputs.logits, dim=-1)
print(f"Probabilities: {probs}")

predicted_class_id = probs.argmax(dim=-1).item()
predicted_label = model.config.id2label[predicted_class_id]
predicted_score = probs[0, predicted_class_id].item()

print(
    f"\nFinal result: {{'label': '{predicted_label}', 'score': "
    f"{predicted_score:.4f}}}"
)
\end{lstlisting}



\textbf{잠깐 --- \inlinecode{transformers} 안에는 softmax 함수가 없나요?}

없습니다. Hugging Face는 \emph{모델·토크나이저·학습 루프} 를 제공하고, \textbf{수치 연산은 PyTorch에 위임}합니다 (혹은 TensorFlow / JAX). 그래서 후처리(softmax, argmax, log 등)는 \inlinecode{torch.*} 를 직접 부릅니다.

PyTorch 안에는 같은 softmax를 표현하는 세 가지 형태가 있습니다 --- 결과는 모두 동일하고 \emph{어디서 호출하느냐} 만 다릅니다.



\Needspace{23\baselineskip}
\begin{lstlisting}
import torch.nn.functional as F

p1 = torch.softmax(outputs.logits, dim=-1)

p2 = F.softmax(outputs.logits, dim=-1)

softmax_module = torch.nn.Softmax(dim=-1)
p3 = softmax_module(outputs.logits)

print(f"Do all three forms agree?")
print(
    f"  torch.softmax vs F.softmax:    "
    f"{torch.allclose(p1, p2)}"
)
print(
    f"  torch.softmax vs nn.Softmax:   "
    f"{torch.allclose(p1, p3)}"
)
print(f"\n  values: {p1}")
\end{lstlisting}



\textbf{보너스 --- 학습 코드에서 자주 보이는 \inlinecode{log\_softmax}}

수치적 안정성 때문에 학습 시에는 \inlinecode{softmax\ →\ log} 두 단계 대신 \textbf{\inlinecode{log\_softmax} 한 번} 으로 묶는 게 표준입니다 (\inlinecode{CrossEntropyLoss} 가 내부에서 이렇게 함).

\Needspace{7\baselineskip}
\begin{lstlisting}
log_probs_unstable = torch.log(torch.softmax(logits, dim=-1))

log_probs_stable = F.log_softmax(logits, dim=-1)
\end{lstlisting}

추론 시에는 단순히 \inlinecode{softmax} 가 명확합니다 --- 확률이 직접 필요하니까요. \textbf{학습 시} \inlinecode{CrossEntropyLoss(logits,\ target)} 는 내부적으로 logit에 \inlinecode{log\_softmax} 를 적용하고 NLL을 더하므로, \emph{softmax를 직접 부를 일이 없습니다} (9장 이후 자주 등장).



\Needspace{9\baselineskip}
\begin{lstlisting}
print(f"pipeline result:    {classifier(text)}")
print(
    f"manual 4-step:      [{{'label': '{predicted_label}', 'score': "
    f"{predicted_score:.4f}}}]"
)
\end{lstlisting}



\section{보너스: 토크나이저마다 어휘가 다르다}

지금까지는 DistilBERT의 WordPiece 토크나이저 \emph{하나} 만 봤습니다. 그런데 모델이 바뀌면 토크나이저도 바뀌고, \textbf{같은 문장이 완전히 다른 토큰 리스트로 쪼개집니다} --- 어휘 사전이 사전학습 단계에서 따로 만들어졌기 때문이에요.

세 가지 대표 토크나이저를 나란히 비교합니다.

\begin{booktable}{7장 보너스: 토크나이저마다 어휘가 다르다 표}{tab:ch07-06}
\begin{adjustbox}{max width=\textwidth}
\begin{tabular}{@{}llll@{}}
\toprule
모델 & 알고리즘 & 어휘 수 & 대소문자 \\
\midrule

\inlinecode{distilbert-base-uncased} & \textbf{WordPiece} & 30,522 & 모두 소문자로 \\
\inlinecode{bert-base-cased} & \textbf{WordPiece} & 28,996 & 대소문자 유지 \\
\inlinecode{gpt2} & \textbf{BPE} (Byte Pair Encoding) & 50,257 & 대소문자 유지 \\
\bottomrule
\end{tabular}
\end{adjustbox}
\end{booktable}

WordPiece와 BPE는 둘 다 \emph{서브워드 알고리즘} 이지만 학습·표기 방식이 달라서 토큰 모양이 시각적으로도 구분됩니다 --- \inlinecode{\#\#} 접두사 vs \inlinecode{Ġ} (공백) 접두사.



\Needspace{17\baselineskip}
\begin{lstlisting}
tokenizer_specs = {
    "distilbert-base-uncased": AutoTokenizer.from_pretrained("distilbert-base-uncased"),
    "bert-base-cased":         AutoTokenizer.from_pretrained("bert-base-cased"),
    "gpt2":                    AutoTokenizer.from_pretrained("gpt2"),
}

print(f"{'model':>28}  {'vocab_size':>10}  {'class':>32}")
print("-" * 76)
for name, tok in tokenizer_specs.items():
    print(
        f"{name:>28}  {tok.vocab_size:>10,}  "
        f"{type(tok).__name__:>32}"
    )
\end{lstlisting}

\begin{codeRead}
\textbf{1--5행}의 \inlinecode{tokenizer\_specs = ...}에서는 이후 분석에서 사용할 값을 준비합니다. \textbf{9행}의 \inlinecode{for name, tok in tokenize...}에서는 앞 단계에서 만든 값을 바탕으로 다음 계산을 수행합니다.
\end{codeRead}



\Needspace{18\baselineskip}
\begin{lstlisting}
sample_sentences = [
    "I love using Hugging Face!",
    "Tokenization is fascinating.",
]

for sent in sample_sentences:
    print(f"Input: {sent!r}")
    for name, tok in tokenizer_specs.items():
        tokens = tok.tokenize(sent)
        print(
            f"  {name:>28}  ({len(tokens)} tokens) "
            f"{tokens}"
        )
    print()
\end{lstlisting}

\begin{codeRead}
\textbf{1--4행}의 \inlinecode{sample\_sentences = ...}에서는 이후 분석에서 사용할 값을 준비합니다. \textbf{6행}의 \inlinecode{for sent in sample\_senten...}에서는 앞 단계에서 만든 값을 바탕으로 다음 계산을 수행합니다. \textbf{8행}의 \inlinecode{for name, tok in tokenize...}에서는 앞 단계에서 만든 값을 바탕으로 다음 계산을 수행합니다. \textbf{9행}의 \inlinecode{tokens = ...}에서는 이후 분석에서 사용할 값을 준비합니다.
\end{codeRead}



\subsection{특수 토큰(special token)이란}

\texttt{{[}CLS{]}}, \texttt{{[}SEP{]}} 같은 토큰은 \emph{문장 텍스트} 가 아니라 \textbf{모델에 신호를 주기 위해 사전학습 단계에서 정해진 약속} 입니다. 어휘 사전에 별도 ID로 들어 있고, 토크나이저가 입력에 자동으로 붙입니다.

\begin{booktable}{7장 특수 토큰(special token)이란 표}{tab:ch07-07}
\begin{adjustbox}{max width=\textwidth}
\begin{tabular}{@{}llll@{}}
\toprule
토큰 & 풀이름 & 위치 & 역할 \\
\midrule

\texttt{{[}CLS{]}} & Classification & 모든 입력 \emph{맨 앞} & 분류 헤드는 \emph{이 위치} 의 hidden state를 사용. attention을 통해 전체 문장 정보가 {[}CLS{]}로 모이도록 학습됨. \\
\texttt{{[}SEP{]}} & Separator & 문장 끝, 두 문장 사이 & 한 문장 입력엔 \texttt{{[}CLS{]}\ ...\ {[}SEP{]}}. 두 문장이면 \texttt{{[}CLS{]}\ A\ {[}SEP{]}\ B\ {[}SEP{]}} (NSP·QA·NLI 등). \\
\texttt{{[}PAD{]}} & Padding & 짧은 문장 끝 & 배치 안 문장 길이를 맞추는 더미 토큰. \textbf{\inlinecode{attention\_mask=0}} 으로 표시해 모델이 무시. \\
\texttt{{[}UNK{]}} & Unknown & 어디든 & 어휘 사전에 없는 토큰. WordPiece는 거의 항상 더 작은 서브워드로 쪼개므로 실제 출현은 드뭄. \\
\texttt{{[}MASK{]}} & Mask & 사전학습 시 입력 일부 & BERT 사전학습의 \emph{Masked LM} --- 입력 토큰 15\%를 \texttt{{[}MASK{]}} 로 가리고 모델이 맞추도록. 추론 시엔 거의 안 등장(fill-mask 데모 제외). \\
\bottomrule
\end{tabular}
\end{adjustbox}
\end{booktable}

\textbf{autoregressive 모델 (GPT-2)} 은 \texttt{{[}CLS{]}/{[}SEP{]}} 가 없습니다 --- 다음 토큰을 \emph{순서대로} 예측하는 구조라 문장 시작/끝 마커가 별도로 필요 없고, \texttt{\textless{}\textbar{}endoftext\textbar{}\textgreater{}} 라는 단일 토큰이 BOS/EOS 역할을 겸합니다.

이 약속은 \emph{모델별로 다릅니다}. RoBERTa는 \texttt{\textless{}s\textgreater{}}, \texttt{\textless{}/s\textgreater{}} 를, T5는 \texttt{\textless{}pad\textgreater{}}, \texttt{\textless{}extra\_id\_0\textgreater{}} 등을 씁니다 --- \inlinecode{tokenizer.special\_tokens\_map} 으로 한 번에 확인 가능.



\Needspace{25\baselineskip}
\begin{lstlisting}
print(
    f"{'model':>28}  {'BOS/CLS':>16}  {'EOS/SEP':>16}  {'PAD':>10}  "
    f"{'UNK':>10}"
)
print("-" * 90)
for name, tok in tokenizer_specs.items():
    cls = tok.cls_token if tok.cls_token else (tok.bos_token or "—")
    sep = tok.sep_token if tok.sep_token else (tok.eos_token or "—")
    pad = tok.pad_token or "—"
    unk = tok.unk_token or "—"
    print(
        f"{name:>28}  {cls:>16}  {sep:>16}  {pad:>10}  "
        f"{unk:>10}"
    )

print()
for name, tok in tokenizer_specs.items():
    print(
        f"{name}.special_tokens_map = "
        f"{tok.special_tokens_map}"
    )
\end{lstlisting}



\textbf{관찰 포인트}

\begin{itemize}
\tightlist
\item
  \textbf{\inlinecode{\#\#} 접두사 (WordPiece)}: DistilBERT·BERT는 단어 중간 서브워드를 \inlinecode{\#\#xxx} 로 표시. 예: \texttt{tokenization\ →\ {[}\textquotesingle{}token\textquotesingle{},\ \textquotesingle{}\#\#ization\textquotesingle{}{]}}. 이전 토큰의 \emph{연속} 이라는 신호.
\item
  \textbf{\inlinecode{Ġ} 접두사 (BPE)}: GPT-2는 토큰 앞에 공백이 있었는지를 \inlinecode{Ġ} (Latin small letter G with stroke) 로 표시. 예: \inlinecode{Ġlove} 는 “love 앞에 공백이 있었다”. 토큰화/디코딩이 정확히 역연산이 되도록 하는 표기.
\item
  \textbf{대소문자}: \inlinecode{bert-base-cased} 는 \inlinecode{Hugging\ Face} 의 \inlinecode{H}, \inlinecode{F} 를 그대로 보존. \inlinecode{distilbert-base-uncased} 는 모두 소문자. 이름·고유명사 처리에서 차이가 큽니다.
\item
  \textbf{vocab 크기}: GPT-2가 50K 로 가장 큼. BPE는 영어 외 다양한 토큰(드문 조합·바이트 단위)도 어휘에 포함하기 때문. WordPiece는 영어 중심이라 30K로 충분.
\item
  \textbf{특수 토큰}: BERT 계열은 \texttt{{[}CLS{]}}, \texttt{{[}SEP{]}}, \texttt{{[}PAD{]}}, \texttt{{[}UNK{]}} 가 모두 정의되지만 GPT-2는 \texttt{{[}CLS{]}/{[}SEP{]}} 가 없습니다 (autoregressive 모델은 문장 시작/끝 마커를 따로 안 둠 --- \texttt{\textless{}\textbar{}endoftext\textbar{}\textgreater{}} 하나가 BOS/EOS 역할). PAD도 없어 추가 설정이 필요한 경우가 흔함.
\end{itemize}

\textbf{왜 같은 문장이 다른 토큰 시퀀스가 되나?} 어휘 사전이 \emph{사전학습 데이터} 에 따라 만들어집니다.

\begin{itemize}
\tightlist
\item
  BERT는 BookCorpus + Wikipedia로 학습됐고, 영어 중심 어휘.
\item
  GPT-2는 더 다양한 웹 텍스트(Reddit 등)로 학습됐고 BPE라 어휘가 더 풍부.
\item
  한국어 BERT(\inlinecode{klue/bert-base}, 14장)는 한국어 코퍼스로 다시 학습돼 한국어 어휘를 보유 --- 같은 문장 \inlinecode{"안녕"} 도 영어 BERT면 \texttt{{[}UNK{]}} 또는 글자 단위로 쪼개지지만 한국어 BERT엔 한 토큰으로 들어갑니다.
\end{itemize}

\textbf{실무 함의}: 모델을 갈아 끼울 때 토크나이저도 \emph{반드시 짝} 으로 바꿔야 합니다. \inlinecode{AutoTokenizer.from\_pretrained(model\_name)} 의 model\_name 이 모델 자체와 일치해야 하는 이유 --- 학습 때 본 어휘와 추론 때 입력 어휘가 같아야 모델이 의미를 이해합니다.



\section{\texorpdfstring{보너스: \inlinecode{model.config} 안에 뭐가 있나}{보너스: model.config 안에 뭐가 있나}}

위에서 \inlinecode{model.config.id2label} 로 라벨 이름을 알아냈습니다. \inlinecode{config} 객체에는 모델의 \emph{설계도} 가 모두 들어있어서, 모델을 받아왔을 때 가장 먼저 들여다보면 좋은 곳입니다.

분류 작업에서 자주 쓰는 속성들을 한 번에 출력합니다.



\Needspace{26\baselineskip}
\begin{lstlisting}
cfg = model.config
n_params, size_mb = model_size_summary(model)

print(f"Model name/path:          {cfg._name_or_path}")
print(f"Model type:               {cfg.model_type}")
print(
    f"Parameters:               {n_params:,}  ("
    f"{n_params/1e6:.1f} M)"
)
print(
    f"fp32 size:                "
    f"{size_mb:.1f} MB  (= params x 4 bytes)"
)
print(
    f"hidden_size:             "
    f"{cfg.hidden_size}        (BERT-base/DistilBERT: 768)"
)
print(
    f"vocab_size:              "
    f"{cfg.vocab_size:,}     (matches tokenizer vocab)"
)
print(
    f"max_position_embeddings: "
    f"{cfg.max_position_embeddings}  (input length cap)"
)
print(
    f"num_labels:              "
    f"{cfg.num_labels}          (classification head dim)"
)
print(f"id2label:                {cfg.id2label}")
print(f"label2id:                {cfg.label2id}")
print(
    f"problem_type:            "
    f"{cfg.problem_type!r}    (None → auto-inferred from num_labels)"
)
\end{lstlisting}



\begin{quote}
📒 \textbf{더 깊이 보고 싶다면 --- 부록 노트북}

\href{./appendix_model_config.ipynb}{\inlinecode{appendix\_model\_config.ipynb}} 에서 다음을 다룹니다:

\begin{itemize}
\tightlist
\item
  \inlinecode{PretrainedConfig} 의 정체와 클래스 계층 (BertConfig / GPT2Config / T5Config / ViTConfig \ldots)
\item
  \inlinecode{AutoConfig.from\_pretrained} 로 \emph{가중치 없이} config만 로드
\item
  5종 모델(bert / distilbert / gpt2 / t5 / roberta) config를 한 표에 비교 + ViT(비전) 사례
\item
  분류 헤드 갈아끼우는 \inlinecode{from\_pretrained} 인자 패턴 (\inlinecode{num\_labels}, \inlinecode{problem\_type})
\item
  공식 문서 링크: \url{https://huggingface.co/docs/transformers/en/main_classes/configuration}
\end{itemize}

Colab으로 바로: \href{https://colab.research.google.com/github/yoon-gu/neuqes-101/blob/master/07_bert_pipeline/appendix_model_config.ipynb}{Open}. 본 장 흐름과 별개라 시간 될 때 보시면 됩니다.
\end{quote}

\textbf{자주 쓰는 속성 한눈에 보기}

\begin{booktable}{7장 보너스: model.config 안에 뭐가 있나 표}{tab:ch07-08}
\begin{adjustbox}{max width=\textwidth}
\begin{tabular}{@{}lll@{}}
\toprule
속성·호출 & 의미 & 자주 쓰는 곳 \\
\midrule

\inlinecode{model.config.\_name\_or\_path} & 모델 식별자 (Hugging Face Hub repo 또는 로컬 경로) & 어떤 모델인지 빠르게 확인 \\
\inlinecode{model.config.model\_type} & 모델 아키텍처 종류 (\inlinecode{bert}, \inlinecode{distilbert}, \inlinecode{gpt2}, ...) & 분기 처리 \\
\inlinecode{sum(p.numel()\ for\ p\ in\ model.parameters())} & \textbf{파라미터 총 개수} (config 속성은 아니지만 항상 같이 봄) & VRAM 사용량 추정, 모델 비교 \\
\inlinecode{model.config.hidden\_size} & hidden state 차원 (예: 768 / 1024) & 분류 헤드를 직접 만들 때 \\
\inlinecode{model.config.vocab\_size} & 어휘 수 (토크나이저와 일치해야 함) & 토크나이저 호환 검증 \\
\inlinecode{model.config.max\_position\_embeddings} & 입력 토큰 수 상한 & \inlinecode{truncation=True,\ max\_length=...} 결정 \\
\inlinecode{model.config.num\_labels} & 분류 헤드 출력 클래스 수 & 모델 로드 시 명시: \inlinecode{num\_labels=5} \\
\inlinecode{model.config.id2label} / \inlinecode{label2id} & 클래스 인덱스 ↔ 이름 매핑 & 추론 결과 해석, 학습 후 모델 카드 친절도 \\
\inlinecode{model.config.problem\_type} & \inlinecode{"regression"} / \inlinecode{"single\_label\_classification"} / \inlinecode{"multi\_label\_classification"} --- \inlinecode{Trainer} 가 자동 loss 결정 & 9·11장·12에서 명시적으로 사용 \\
\bottomrule
\end{tabular}
\end{adjustbox}
\end{booktable}

\textbf{실무 패턴}: 새 모델을 받자마자 \inlinecode{print(model.config)} 또는 \inlinecode{cfg.to\_dict()} 로 내용을 먼저 본다 → 입력/출력 가정을 확인하고 토크나이저·\inlinecode{Trainer} 설정과 일치시킴.

\Needspace{6\baselineskip}
\begin{lstlisting}
print(model.config)
print(model.config.to_dict())
\end{lstlisting}



\section{이 장에 등장한 라이브러리·함수}

\subsection{\texorpdfstring{\inlinecode{transformers}}{transformers}}

\begin{booktable}{7장 transformers 표}{tab:ch07-09}
\begin{adjustbox}{max width=\textwidth}
\begin{tabular}{@{}lll@{}}
\toprule
이름 & 한 줄 설명 & 다음 장에서 \\
\midrule

\inlinecode{transformers.pipeline} & 추론 원스톱 함수 (3단계 묶음) & 학습된 모델을 \inlinecode{pipeline}으로 감싸 사용 (9장 이후) \\
\inlinecode{transformers.AutoTokenizer} & 모델에 맞는 토크나이저 자동 로드 & 8장에서 옵션 깊게 보기 \\
\inlinecode{transformers.AutoModelForSequenceClassification} & 분류 헤드가 붙은 모델 자동 로드 & 9장부터 직접 파인튜닝 \\
\bottomrule
\end{tabular}
\end{adjustbox}
\end{booktable}

\subsection{\texorpdfstring{\inlinecode{torch} 의 후처리·연산 함수}{torch 의 후처리·연산 함수}}

\begin{booktable}{7장 torch 의 후처리·연산 함수 표}{tab:ch07-10}
\begin{adjustbox}{max width=\textwidth}
\begin{tabular}{@{}lll@{}}
\toprule
형태 & 호출 예시 & 언제 \\
\midrule

\inlinecode{torch.softmax(x,\ dim=-1)} & 텐서에 직접 & 추론 후처리 (이 장처럼) \\
\inlinecode{torch.nn.functional.softmax} (보통 \inlinecode{F.softmax}) & \inlinecode{F.softmax(x,\ dim=-1)} & 함수형, PyTorch 코드에 가장 흔함 \\
\inlinecode{torch.nn.Softmax(dim=-1)} & layer로 모델 안에 박을 때 & 잘 안 씀 (보통 logits 그대로 두고 loss가 처리) \\
\inlinecode{F.log\_softmax}, \inlinecode{torch.log\_softmax} & softmax + log 한 번에 (수치 안정) & 학습 loss 직접 구현 시 \\
\inlinecode{torch.argmax(x,\ dim=-1)} 또는 \inlinecode{x.argmax(dim=-1)} & 가장 큰 값의 인덱스 & 분류 예측 인덱스 추출 \\
\inlinecode{torch.no\_grad()} (context manager) & 추론 중 gradient 비활성 & 메모리·속도 절약 \\
\bottomrule
\end{tabular}
\end{adjustbox}
\end{booktable}

\begin{quote}
\textbf{요점}: HuggingFace는 softmax·argmax 같은 \emph{수치 함수} 를 따로 제공하지 않습니다. 모두 PyTorch에서 직접 호출합니다 (TensorFlow 백엔드를 쓰면 \inlinecode{tf.nn.softmax} 같은 식). \inlinecode{Trainer} 가 학습 loss 안에서 softmax를 자동 처리하므로, 학습 코드에서는 직접 부를 일이 거의 없고 \emph{추론 후처리} 에서 등장하는 게 보통.
\end{quote}

\subsection{\texorpdfstring{\inlinecode{model.config} 의 자주 쓰는 속성}{model.config 의 자주 쓰는 속성}}

\begin{booktable}{7장 model.config 의 자주 쓰는 속성 표}{tab:ch07-11}
\begin{adjustbox}{max width=\textwidth}
\begin{tabular}{@{}ll@{}}
\toprule
속성 & 용도 \\
\midrule

\inlinecode{id2label}, \inlinecode{label2id} & 클래스 인덱스 ↔ 이름 매핑 \\
\inlinecode{num\_labels} & 분류 헤드 출력 차원 \\
\inlinecode{hidden\_size}, \inlinecode{vocab\_size}, \inlinecode{max\_position\_embeddings} & 모델 구조 파라미터 \\
\inlinecode{model\_type}, \inlinecode{\_name\_or\_path} & 모델 정체성 \\
\inlinecode{problem\_type} & \inlinecode{Trainer} 자동 loss 선택 (\inlinecode{regression} / \inlinecode{single\_label\_classification} / \inlinecode{multi\_label\_classification}) \\
\bottomrule
\end{tabular}
\end{adjustbox}
\end{booktable}

\subsection{\texorpdfstring{\inlinecode{torch} 자체}{torch 자체}}

\begin{booktable}{7장 torch 자체 표}{tab:ch07-12}
\begin{adjustbox}{max width=\textwidth}
\begin{tabular}{@{}lll@{}}
\toprule
이름 & 한 줄 설명 & 다음 장에서 \\
\midrule

\inlinecode{torch} & PyTorch --- 텐서 연산, 자동 미분, GPU 연결 & 계속 사용 (특히 9장 학습부터) \\
\bottomrule
\end{tabular}
\end{adjustbox}
\end{booktable}



\section{체크포인트 질문}

\begin{enumerate}
\def\labelenumi{\arabic{enumi}.}
\tightlist
\item
  \inlinecode{pipeline("sentiment-analysis")} 을 직접 풀어 쓰면 어떤 단계가 되는지 4단계로 설명할 수 있나요?
\item
  \inlinecode{input\_ids} 와 \inlinecode{attention\_mask} 는 각각 무슨 역할인가요?
\item
  모델 출력의 \inlinecode{logits} 는 왜 그대로 쓰지 않고 \inlinecode{softmax} 를 거치나요? (4장에서 본 동등성 떠올리기)
\item
  \inlinecode{AutoTokenizer.from\_pretrained(...)} 를 쓸 때와 \inlinecode{BertTokenizer.from\_pretrained(...)} 를 직접 쓸 때의 차이는 무엇인가요?
\end{enumerate}



\section{FAQ}

\begin{faqBox}{Q1. (실무) \inlinecode{pipeline} 이 처음 실행될 때 너무 느린데 정상인가요?}

네, 정상입니다. 첫 실행 시 다음이 한꺼번에 일어납니다.

\begin{enumerate}
\def\labelenumi{\arabic{enumi}.}
\tightlist
\item
  모델 가중치 다운로드 (약 250MB for DistilBERT)
\item
  토크나이저 사전 다운로드
\item
  모델 PyTorch 로드 + (있다면) GPU 이동
\end{enumerate}

두 번째부터는 캐시(\texttt{\textasciitilde{}/.cache/huggingface/})에서 읽으므로 즉시 실행됩니다. \textbf{Colab 세션이 끊기면 캐시도 사라져} 다시 다운로드합니다 --- 자주 쓴다면 Google Drive를 마운트해서 \inlinecode{HF\_HOME} 환경변수를 Drive로 지정하면 보존됩니다.

\Needspace{8\baselineskip}
\begin{lstlisting}
import os
from google.colab import drive
drive.mount("/content/drive")
os.environ["HF_HOME"] = "/content/drive/MyDrive/hf_cache"
\end{lstlisting}

\end{faqBox}

\begin{faqBox}{Q2. (이론) 왜 BERT 토크나이저는 단어를 \inlinecode{\#\#} 조각으로 쪼개나요?}

이게 \textbf{WordPiece} 의 핵심입니다 --- 자주 등장하는 부분 문자열을 하나의 토큰으로 두고, 새로운 단어는 \emph{서브워드} 들의 조합으로 표현합니다.

\Needspace{9\baselineskip}
\begin{lstlisting}
tokenizer.tokenize("Tokenization")
# ['token', '##ization']
tokenizer.tokenize("antidisestablishmentarianism")
# → ['anti', '##dis', '##est', '##ab', '##lish', 
# '##ment', '##arian', '##ism']
\end{lstlisting}

\inlinecode{\#\#} 은 “이 토큰은 \emph{이전 토큰의 연속}”이라는 표시입니다.

\textbf{왜 이렇게?}

\begin{itemize}
\tightlist
\item
  \textbf{OOV 해결}: TF-IDF는 학습 어휘에 없는 단어를 \emph{무시} 했지만, WordPiece는 항상 더 작은 서브워드로 쪼개 표현 가능 (이론적으로 OOV 없음, 최악의 경우 글자 단위까지).
\item
  \textbf{어휘 수 관리}: 영어에는 수백만 단어가 있지만 BERT는 30,522개 토큰만으로 모두 표현. 형태 변형(\inlinecode{-ing}, \inlinecode{-ed})도 일관되게 처리.
\item
  \textbf{희귀 단어 일반화}: \inlinecode{unhappiness} 를 \inlinecode{un\ +\ happi\ +\ ness} 로 보면, 모델이 \inlinecode{happi} 를 알고 있으면 처음 보는 단어라도 의미 추론 가능.
\end{itemize}

\end{faqBox}

\begin{faqBox}{Q3. (실무) \inlinecode{pipeline} 결과의 \inlinecode{LABEL\_0}, \inlinecode{LABEL\_1} 이 무슨 의미인지 어떻게 알아내나요?}

\inlinecode{model.config.id2label} 을 확인하면 됩니다.

\Needspace{6\baselineskip}
\begin{lstlisting}
print(classifier.model.config.id2label)
\end{lstlisting}

\textbf{모델마다 다릅니다}. 일부 모델은 \inlinecode{LABEL\_0}, \inlinecode{LABEL\_1} 처럼 \emph{추상적인} 이름만 붙어 있어 모델 카드(Hugging Face Hub의 모델 페이지) 또는 학습 데이터셋 라벨 정의를 확인해야 합니다. 우리가 9장 이후 직접 학습할 때는 \inlinecode{id2label} 을 명시적으로 설정해 미래의 사용자에게 친절하게 만들 수 있습니다.

\end{faqBox}

\begin{faqBox}{Q4. (이론) \texttt{{[}CLS{]}} 와 \texttt{{[}SEP{]}} 토큰은 왜 필요한가요?}

BERT의 사전학습 구조에서 비롯된 특수 토큰입니다.

\begin{itemize}
\tightlist
\item
  \textbf{\texttt{{[}CLS{]}}} (Classification): 입력 맨 앞에 항상 붙는 토큰. 사전학습 시 \emph{Next Sentence Prediction} 을 위한 자리였고, 분류 작업에서는 이 위치의 hidden state(전체 문장의 표현)를 분류 헤드에 넣습니다. \emph{모든 토큰의 정보가 attention을 통해 {[}CLS{]}로 모이도록} 학습됨.
\item
  \textbf{\texttt{{[}SEP{]}}} (Separator): 문장 끝에 붙거나, 두 문장을 분리. 입력이 한 문장이면 \texttt{{[}CLS{]}\ ...\ {[}SEP{]}} 구조, 두 문장(질문-답변 등)이면 \texttt{{[}CLS{]}\ ...\ {[}SEP{]}\ ...\ {[}SEP{]}}.
\end{itemize}

이 장의 입력 \inlinecode{"I\ love\ using\ Hugging\ Face!"} 의 token ID 첫 값 \inlinecode{101} 이 \texttt{{[}CLS{]}}, 마지막 \inlinecode{102} 가 \texttt{{[}SEP{]}} --- 위 셀에서 확인했죠.

\inlinecode{AutoTokenizer} 가 이 특수 토큰을 자동으로 추가해주므로 우리가 신경 쓸 일은 거의 없지만, \emph{왜 길이가 입력 단어 수보다 2 길게 나오는지} 이해하는 데는 이 두 토큰을 알아야 합니다.

\end{faqBox}

\begin{faqBox}{Q5. (실무) GPU 없이도 \inlinecode{pipeline} 이 돌아가나요?}

네, CPU에서도 작동합니다. 다만 속도 차이가 큽니다.

\begin{booktable}{7장 FAQ 표}{tab:ch07-13}
\begin{adjustbox}{max width=\textwidth}
\begin{tabular}{@{}ll@{}}
\toprule
환경 & DistilBERT 1문장 추론 \\
\midrule

Colab CPU & \textasciitilde80-150ms \\
Colab T4 GPU & \textasciitilde5-15ms \\
큰 BERT 모델 & CPU는 5-10x 더 느림 \\
\bottomrule
\end{tabular}
\end{adjustbox}
\end{booktable}

\textbf{추론} 만 한다면 CPU도 실용적입니다 (예: API 서빙은 보통 GPU지만 간단한 데모는 CPU). \textbf{학습} 은 거의 항상 GPU 필수 --- 9장 이후 학습할 때는 T4 런타임으로 바꿔야 합니다.

\end{faqBox}

\begin{faqBox}{Q6. (실무) \inlinecode{AutoTokenizer} 와 \inlinecode{BertTokenizer} 를 직접 사용하는 것의 차이는?}

기능적으론 거의 같지만 \textbf{코드 일반성} 이 다릅니다.

\Needspace{9\baselineskip}
\begin{lstlisting}
from transformers import BertTokenizer
tokenizer = BertTokenizer.from_pretrained("bert-base-uncased")

from transformers import AutoTokenizer
tokenizer = AutoTokenizer.from_pretrained("bert-base-uncased")
\end{lstlisting}

실무에선 \textbf{거의 항상 \inlinecode{AutoTokenizer}} 를 씁니다. \inlinecode{Auto} 계열은 모델 카드(\inlinecode{config.json}) 에서 어떤 클래스를 써야 할지 자동 추론하므로, 다른 모델로 갈아끼우는 실험이 매우 쉽습니다. 14장의 한국어 BERT(\inlinecode{klue/bert-base}) 도 같은 패턴으로 로드됩니다 --- 코드 한 줄도 안 바꾸고요.
\end{faqBox}



\section{오류 실험 (선택)}

다음 코드를 돌려보고 에러 메시지를 읽어보세요. 어떤 인자가 빠졌을까요?

\Needspace{6\baselineskip}
\begin{lstlisting}
inputs_bad = tokenizer("I love HF!")
outputs_bad = model(**inputs_bad)
\end{lstlisting}

힌트: 모델은 PyTorch 텐서를 기대하는데, 토크나이저가 기본값으로 무엇을 반환할까요? \inlinecode{tokenizer("...")} 의 기본 반환 형식과 \inlinecode{tokenizer("...",\ return\_tensors="pt")} 의 차이를 출력 비교해보세요.



\begin{previewBox}{미리보기: 다음 장}

\textbf{8장. 토크나이저 옵션과 데이터셋 (Tokenizer \& Datasets)}

\begin{itemize}
\tightlist
\item
  서브워드 토큰화 옵션 --- \inlinecode{padding}, \inlinecode{truncation}, \inlinecode{max\_length} 의 의미
\item
  \inlinecode{datasets} 라이브러리: \inlinecode{load\_dataset}, \inlinecode{map}, \inlinecode{filter}, Apache Arrow 메모리 효율
\item
  DataLoader 변환 --- 9장 학습 코드의 입력 준비 단계
\item
  \textbf{여전히 학습 없음} (추론 데이터 파이프라인까지만)
\end{itemize}
\end{previewBox}